%********************************************************************
% Appendix
%*******************************************************
% % If problems with the headers: get headings in appendix etc. right
\markboth{\spacedlowsmallcaps{Appendix}}{\spacedlowsmallcaps{Appendix}}
\chapter{\prototype{} Example Questions and Answers}
\label{appendix:ch:example_questions}

This chapter contains logs of example prompts and the system's responses. It serves as a showcase of how \prototype{} works. The examples contain the raw responses, returned by \prototype{} in Markdown and Cypher. The examples span three datasources: the examination office's dataset, the textbook information, and the SpeedParcel example dataset.

%%%%%%%%%%%%%
% WHAT DOES HIS DO
\section{What does HIS do?}
\label{appendix:code:whatDoesHisDo}
A basic question about the examination office's dataset, which \prototype{} is able to answer based on the contents of the knowledge graph.

\begin{lstlisting}[style=cypherstyle, label={lst:qWhatIsHis},caption={Asking \prototype{} what the application HIS is.}]
User: What does HIS do?

Agent:
## Summary

Based on the findings, "HIS-POS" is an application component (and also categorized as an ArchiElement) that supports several academic and administrative functions.

## Details

HIS-POS enables or realizes the following application functions:

- Creating and editing module exams
- Viewing and searching student data
- Managing exam booking and grades (including entering, changing, and viewing grades)
- Booking EN ("Ergebnis notwendig") and NB ("Nicht bestanden") statuses

No associated business functions are explicitly documented in this dataset, but the supported application functions indicate that HIS-POS is likely integral to exam and grade management within an academic institution.

\end{lstlisting}

\begin{lstlisting}[style=cypherstyle, label={lst:cWhatIsHis},caption={The Cypher used to retrieve information on the HIS application from the examination office's knowledge graph.}]
MATCH (n)
WHERE (n:ApplicationComponent OR n:ApplicationFunction OR n:ArchiElement OR n:BusinessFunction)
  AND n.name IS NOT NULL
  AND toUpper(n.name) CONTAINS 'HIS'
OPTIONAL MATCH (n)-[:REALIZATION]->(af:ApplicationFunction)
OPTIONAL MATCH (n)-[:SERVING]->(bf:BusinessFunction)
WITH n, collect(DISTINCT af) AS afs, collect(DISTINCT bf) AS bfs
RETURN n.name AS his, labels(n) AS labels, n.documentation AS documentation, [a IN afs | a.name] AS realizes_application_functions, [b IN bfs | {name: b.name, documentation: b.documentation}] AS serves_business_functions
ORDER BY his
\end{lstlisting}




%%%%%%%%%%%%%
% WHAT IS AN ORPHAN
\section{What is an Orphan?}
\label{appendix:code:whatIsAnOrphan}
A textbook question that generated many insights on how to interact with \prototype{} to get the best results. Prepending the question with the hint to stay within the realm of EAM triggered the LLM to explicitly query the information from the textbook, rather than returning a generic answer from the LLM's pretrained data.

\begin{lstlisting}[style=cypherstyle, label={lst:qWhatIsAnOrphan},caption={Asking \prototype{} what an orphan is in the realm of EAM.}]
User: In Enterprise architecture management, what is an orphan?
Agent:
## Summary

In enterprise architecture, an "orphan" refers to an application that does not support any business capability and therefore has no business value.

## Details

- Orphan applications are disconnected from business needs and are typically candidates for rationalization or retirement.
- Identifying orphans is a key activity when analyzing the application landscape to ensure alignment between IT assets and business value.
- In business support matrices, orphans are highlighted as not contributing to any required capabilities.

\end{lstlisting}

\begin{lstlisting}[style=cypherstyle, label={lst:cWhatIsAnOrphan},caption={The Cypher used to query textbook information containing the keyword "orphan".}]
CALL db.index.fulltext.queryNodes("kg_fulltext_main", "orphan OR \"orphan node\" OR orphaned OR \"what is an orphan\"", {limit: 10}) YIELD node, score
WITH node, score, coalesce(node.definition, node.title, node.text, node.table_summary) AS content
OPTIONAL MATCH (sec:Section)-[:HAS_DOCUMENT]->(doc:Document)
WHERE sec = node
RETURN labels(node) AS labels, content, doc.url AS source, score
ORDER BY score DESC
\end{lstlisting}

%%%%%%%%%%%%%
% WHICH APPLICATION SUPPORTS THE MOST CAPABILITIES
\section{Which applications supports the most capabilities?}
\label{appendix:code:whatIsAnOrphan}
A question regarding the SpeedParcel dataset, which shows how \prototype{} is able to fit the Cypher to the dataset and fulfill the requested task.

\begin{lstlisting}[style=cypherstyle, label={lst:qWhichApplication},caption={Asking \prototype{} which application in SpeedParcel's dataset supports the most capabilities.}]
User: Which application supports the most capabilities?
Agent:
## Summary

The application supporting the most capabilities is **Power BI**, with support for **19 distinct capabilities**.

## Details

- **Power BI** is the application that currently has the broadest capability coverage in your portfolio.
- This suggests it plays a significant role in enabling business functions and may be a key candidate for further investment or rationalization analysis.
\end{lstlisting}

\begin{lstlisting}[style=cypherstyle, label={lst:cWhichApplication},caption={The Cypher used to count the amount of capabilities an application supports".}]
MATCH (a:Application)
OPTIONAL MATCH (a)-[:SUPPORTS]->(c:Capability)
WITH a, count(DISTINCT c) AS capabilityCount
ORDER BY capabilityCount DESC, a.name ASC
LIMIT 1
RETURN a.name AS application, capabilityCount
\end{lstlisting}